% sections/05_results.tex — Section V (Experimental Results)
% Sources:
%   notes/paper_section_VA_draft.md  → §V.A (Day 5)
%   notes/paper_section_VB_draft.md  → §V.B (Day 5)
%   §V.C/§V.D placeholders for Day 6 / Day 14+
% Converted Day 5 2026-05-13.

\section{Experimental Results}
\label{sec:results}

\subsection{Diagnostic Measurements: Signal-Flow Across the VLA Pipeline}
\label{sec:results:diagnostic}

\subsubsection{Methodology Overview}
\label{sec:results:diagnostic:methodology}

The three-layer signal-flow diagnostic comprises three Mantel tests on inter-task
distance matrices, applied to different stages of the VLA pipeline
(Table~\ref{tab:diagnostic_layers}).

\begin{table}[t]
\centering
\caption{\textsc{Signal-Flow Diagnostic Layers.}
Each Mantel test compares inter-task distance matrices on a stage-specific
representation against inter-task distance on language embeddings.}
\label{tab:diagnostic_layers}
\small
\begin{tabular}{lll}
\toprule
\textbf{Test} & \textbf{Subject of measurement} & \textbf{Pipeline stage} \\
\midrule
Gate~1 & $h_{\mathrm{backbone}}$ (LLM hidden state) & Pre-training endpoint \\
M3     & Demonstration trajectory $a$            & Supervision signal \\
C1     & Policy latent $z$ (head output)         & Trained policy \\
\bottomrule
\end{tabular}
\end{table}

Each measurement applies the same statistical protocol: (1)~collect representations
$\{X_k\}_{k=1}^{K}$ across $K = 10$ tasks of LIBERO-Spatial; (2)~compute the inter-task
distance matrix $D^X_{ij} = \|X_i - X_j\|$; (3)~compute the matched inter-task distance
matrix on language embeddings $D^{\ell}_{ij} = \|\ell_i - \ell_j\|$; (4)~Mantel test on
$(D^X, D^{\ell})$ with $N_{\text{perm}} = 1000$ permutations to obtain Pearson correlation
$r$ and permutation $p$-value. Distances use Euclidean norm; language embeddings $\ell_i$
are mean-pooled token embeddings from the OpenVLA-OFT tokenizer applied to each task's
instruction prompt.

\subsubsection{Gate~1 --- Backbone Language Preservation}
\label{sec:results:diagnostic:gate1}

\paragraph{Methodology.}
For each task, we tokenize the instruction prompt with the OpenVLA-OFT processor, forward
the resulting tokens through the OFT vision-language model (frozen backbone, no
fine-tuning), and extract the LLM hidden state at position $[-1]$ --- the last text token
before the action prediction segment. This position aggregates causal attention over the
entire visual+text input sequence. Per task, we collect $n_{\text{demos}} = 10$ samples
drawn from the first frame of 10 demonstration episodes; the per-task representation is
the mean over these samples.

\paragraph{Sanity checks.}
Three sanity checks (SC1, SC2, SC3) precede the main Gate~1 measurement:
SC1 verifies position $[-1]$ corresponds to the last text token (id~29901, surface form
\texttt{:}), which aggregates causal attention over the full visual + text sequence;
SC2 verifies that the language embedding pipeline is deterministic across two
independent runs ($L_2$ norm $= 0.234$, max element-wise diff $= 0$); SC3 verifies
replication of the original Gate~1 result ($r = 0.7575$ vs.\ original $r = 0.7423$,
$\Delta r = 0.015$ within sampling noise).

\paragraph{Result.}
\begin{equation}
r(h_{\mathrm{backbone}},\, \ell) \;\approx\; +0.76, \qquad p < 0.001.
\label{eq:gate1_result}
\end{equation}
The squared correlation $\rho_h^2 = 0.578$ implies a Gaussian-equivalent MI surrogate of
$\hat{I}_h \approx 0.43$ nats (qualitative reference only;
Remark~\ref{rem:approx_bound}). The OpenVLA-OFT backbone, prior to any
demonstration-specific fine-tuning of downstream modules, preserves strong correlation
between its hidden state and the language instruction across tasks.

\subsubsection{M3 --- Demonstration-Level Language-Action Correlation}
\label{sec:results:diagnostic:m3}

\paragraph{Methodology.}
For each task we collect $n_{\text{demos}} = 10$ demonstration episodes from
LIBERO-Spatial and compute three trajectory representations: flat (concatenated action
chunk, primary), endpoint (first and last actions), and waypoint (subsampled actions).
Per task, the representation is the mean over 10 episodes.

\paragraph{Result.}
Across the three representations,
\begin{equation}
r(\text{traj}, \ell) \;\in\; [-0.27,\, -0.20], \qquad \rho^2 \le 0.073.
\label{eq:m3_result}
\end{equation}
The maximum-magnitude correlation is on the flat representation: $r = -0.2698$ with
permutation $p = 0.964$. Mean per-task trajectory $L_2$ norms are tightly clustered
(range $[20.376, 20.868]$), ruling out a degenerate explanation where M3's near-zero
correlation arises from trajectory-norm artifacts.

This is the empirical realization of the $\eps$-language-orthogonality condition
(Definition~\ref{def:eps_lang_orthog}) with $\eps$ small enough that the Gaussian MI
surrogate is $\hat{I}_{\dataD}(a; \ell \mid o) \approx 0.04$ nats. The
squared-correlation gap
\begin{equation}
\rho_h^2 \,/\, \rho_{\text{demo}}^2 \;=\; 0.578 \,/\, 0.073 \;\approx\; 8\times
\label{eq:gap_result}
\end{equation}
is the principal effect-size measurement of this paper.

\subsubsection{C1 --- Policy-Level Language Correlation in the Trained Latent}
\label{sec:results:diagnostic:c1}

\paragraph{Methodology.}
We train OTP-Soft V3 (Shortcut Flow Matching head + decoder, full architecture as in
\S\ref{sec:architecture:otp_soft}) on LIBERO-Spatial demonstrations, forward 10
evaluation samples per task through the trained policy, extract the head output $z$, and
compute inter-task distance matrices on $z$. We also compute $\eta^2$ ---
the proportion of variance in each $z$ slot attributable to task identity.

\paragraph{Result.}
\begin{equation}
r(z_{\text{OTP}}, \ell) \;=\; -0.19, \qquad p_{\text{Mantel}} = 0.913.
\label{eq:c1_result}
\end{equation}
The $\eta^2$ analysis reveals structure: median $\eta^2$ across slots is $0.26$, range
$[0.05, 0.96]$, with high-$\eta^2$ slots discriminating tasks via scene geometry (where
the bowl and plate are positioned) rather than via what the instruction says about them.
The task-discriminative variance is concentrated in slots corresponding to non-target
object pose encoding.

The trained policy's latent encodes task-discriminative information, but this information
is visually grounded rather than linguistically driven. Condition~\ref{cond:c2} is
\emph{not} substantively satisfied for the OTP-Soft V3 policy.

\subsubsection{Backbone Capacity Check}
\label{sec:results:diagnostic:capacity}

A potential concern with the Gate~1 result is whether the high backbone correlation
reflects genuine task-discriminative encoding or merely the encoding of task-correlated
visual features available in the prompt. We address this by training a width-doubled
deterministic head on a controlled probe:
\begin{equation}
\bar{F}^a \;=\; 11.30, \qquad \max_a F^a \;=\; 26.90.
\label{eq:test4_result}
\end{equation}
The mean per-dimension $F$-statistic of $11.30$ across 7 action dimensions exceeds the
null threshold by a substantial margin. The M3 result is a property of the supervision
distribution, not of the backbone's representational capacity.

\subsubsection{Summary}
\label{sec:results:diagnostic:summary}

\begin{table}[t]
\centering
\caption{\textsc{Three-Layer Signal-Flow Summary.}}
\label{tab:diagnostic_summary}
\small
\begin{tabular}{lll}
\toprule
\textbf{Stage} & \textbf{Correlation with $\ell$} & \textbf{Interpretation} \\
\midrule
Backbone (pre-training)        & $r = +0.76$  & Language preserved \\
Demonstration (supervision)    & $r = -0.27$  & Language attenuated \\
Policy latent (post-training)  & $r = -0.19$  & Inherits demo attenuation \\
\bottomrule
\end{tabular}
\end{table}

The transition from $r = +0.76$ to $r = -0.27$ between backbone and demonstrations
identifies the supervision data as the locus of information loss, consistent with
Theorem~\ref{thm:bound} bounding policy-level language MI by demonstration-level MI.

\subsection{Paraphrase Invariance of OpenVLA-OFT}
\label{sec:results:paraphrase}

\subsubsection{Experimental Setup}
\label{sec:results:paraphrase:setup}

To assess the standard interpretation that paraphrase invariance evidences ``language
understanding'' in VLA policies, we evaluate OpenVLA-OFT~\cite{kim2025openvla_oft} on
LIBERO-Spatial under controlled instruction rewording. The experimental matrix comprises
10 tasks $\times$ 6 instruction phrasings $\times$ 50 evaluation episodes per
task-phrasing cell, totaling 3000 episode evaluations.

The 6 phrasings per task consist of one ``identity'' instruction (the original
LIBERO-Spatial task description used for training) and 5 semantically-equivalent
paraphrasings. Paraphrasings vary surface lexical form (e.g., ``pick up'' $\to$
``grasp''; ``place'' $\to$ ``put'') while preserving semantic content. Each task-phrasing
cell is evaluated under 50 independent episode rollouts from the LIBERO-Spatial test
split using the published OpenVLA-OFT checkpoint without further fine-tuning. We report
point-estimate task success rate over 50 episodes plus Wilson 95\% confidence intervals,
and aggregate paired identity-vs-paraphrased differences with 95\% CIs across the task
$\times$ phrasing matrix.

\subsubsection{Aggregate Result --- Paraphrase Invariance}
\label{sec:results:paraphrase:aggregate}

Across the full 10-task $\times$ 6-phrasing matrix:
\begin{equation}
\text{Identity SR} \;=\; 94.6\%, \quad \text{Paraphrased SR} \;=\; 93.7\%, \quad
\Delta \;=\; 0.9\,\text{pp}.
\label{eq:paraphrase_aggregate}
\end{equation}

The aggregate identity-vs-paraphrased difference of 0.9 percentage points is small
relative to the per-cell sampling variance (with $n = 50$ episodes per cell, the Wilson
95\% CI half-width is approximately $\pm 9$pp). Under any standard practical-equivalence
threshold ($\pm 2$pp, $\pm 3$pp, $\pm 5$pp), OFT's performance is statistically
indistinguishable across paraphrasings.

We formalize this with the Stage~3 statistical protocol of our pre-registration. The
per-cell BH-FDR-adjusted significance test shows that 0 of 50 task-phrasing cells achieve
$q < 0.05$ for any pairwise identity-vs-paraphrasing comparison (McNemar test with
Benjamini--Hochberg adjustment across 50 cells). All 5 paired identity-vs-paraphrased SR
differences (one per paraphrasing variant, aggregated across tasks) have 95\% confidence
intervals that include the practical-equivalence threshold $\pm 2$pp. By the
pre-registered protocol, the joint conditions for ``visual-grounded SR regime'' (the
policy achieves high SR robust to instruction paraphrasing) are \textbf{supported} by
these observations.

\subsubsection{Task 5 Sensitivity Caveat}
\label{sec:results:paraphrase:task5}

One task in the 10-task LIBERO-Spatial suite (``pick up the black bowl on the stove and
place it on the plate,'' task 5) exhibits an unusual pattern: identity SR is 60\% rather
than the 90\%+ observed on other tasks, with paraphrased SR also showing increased
variance. This appears to be a per-task difficulty artifact unrelated to
language-conditioning behavior --- the same physical task is challenging across all
phrasing variants --- but it warrants explicit sensitivity analysis.

Excluding task 5 from the aggregate produces:
\begin{equation}
\text{Identity SR (9-task)} = 98\%, \quad
\text{Paraphrased SR (9-task)} = 98\%, \quad
\Delta = 0\,\text{pp}.
\end{equation}
The 9-task subset preserves and strengthens the paraphrase-invariance finding. Excluding
task 5 does not change any other conclusion in this section. We retain the full 10-task
analysis as primary and report the 9-task subset as a sensitivity check rather than as a
post-hoc selected analysis.

\subsubsection{Interpretation Under Theorem~\ref{thm:bound}}
\label{sec:results:paraphrase:interpretation}

The paraphrase-invariance observation is conventionally read as evidence that the policy
``understands'' language: a model robust to instruction rewording is interpreted as
having learned semantic content beyond surface form. Theorem~\ref{thm:bound} sharpens
this interpretation. Under the bound $I(\Astar; \ell \mid o) \le I_{\dataD}(a; \ell \mid
o)$ with demonstration-level signal $I_{\dataD}(a; \ell \mid o) \approx 0.04$ nats on
LIBERO-Spatial (\S\ref{sec:results:diagnostic:m3}), OFT's 94.6\%/93.7\%
identity-vs-paraphrased SR result is consistent with two distinct scenarios:

\textbf{(a) Low-magnitude language preservation.} OFT does carry some language
conditioning within the bound, sufficient to discriminate between semantically-different
instructions but not surface-form variants of the same instruction. Paraphrase invariance
reflects genuine semantic conditioning at low MI magnitude.

\textbf{(b) Entirely visual grounding.} OFT carries no language conditioning at the
action level. Paraphrase invariance reflects that the policy ignores language entirely
and acts based on observation alone. With LIBERO-Spatial's tightly task-correlated visual
scenes (each task has a distinctive object layout), observation alone is sufficient for
task discrimination.

The paraphrase-invariance measurement cannot distinguish (a) from (b). Both predict the
observed pattern: high identity SR, paraphrased SR statistically equivalent to identity
SR. The two scenarios differ only in the \emph{mechanism} of task success, which is not
directly probed by the SR metric.

This is the diagnostic methodology contribution made explicit in
\S\ref{sec:intro:contributions} (Contribution~4): the SR + paraphrase-invariance
evaluation regime cannot distinguish between language-conditioned and visually-grounded
high performance. Signal-flow diagnostics (Gate~1, M3, C1) provide complementary
information that locates the dominant pathway as visual-grounded for the OFT case (b),
though without ruling out small residual language preservation (a) at low MI magnitude.

\subsubsection{Implications}
\label{sec:results:paraphrase:implications}

\textbf{1.~Paraphrase invariance is not evidence of language conditioning.} Under
Theorem~\ref{thm:bound}, any policy with demonstration-supervised training on
LIBERO-Spatial-like data is bounded to small policy-level language MI. Both
linguistically-conditioned policies (at low magnitude within the bound) and entirely
visually-grounded policies predict paraphrase invariance.

\textbf{2.~High SR is not evidence of language conditioning.} OpenVLA-OFT's 97\% SR
coexists with the Theorem~\ref{thm:bound} bound: the policy achieves task success via
observation-action mapping; the contribution of language to action selection is bounded
above by $\sim 0.04$ nats. Standard benchmark SR cannot establish language-driven task
completion.

\textbf{3.~Signal-flow diagnostics are complementary, not replacements.} Gate~1, M3, and
C1 measure quantities orthogonal to task SR. They localize where language information is
preserved, lost, or recovered across the VLA pipeline. We recommend their adoption
alongside, not in place of, standard SR + paraphrase-invariance evaluations.

\subsection{Correlation Gap, OFT Contrast, and CFM Decoder Collapse}
\label{sec:results:correlation}
\label{sec:results:oft_contrast}
\label{sec:results:cfm_collapse}

% Source: To be drafted Day 6+ from §V.A and §V.B integrated narrative.

[To be added Day~6: (a) detailed correlation-gap calculation tying together
\S\ref{sec:results:diagnostic:m3} and \S\ref{sec:results:paraphrase}; (b) OFT contrast
formal exposition --- both within-bound scenarios (low-magnitude language preservation
vs.\ entirely visual grounding) not distinguishable by SR alone; (c) CFM decoder
sample-quality collapse --- 7/7 action dimensions fail per-dimension threshold (overall
L1/std $= 1.88$) under language-orthogonal supervision, oracle ablation (replacing head
output with ground-truth trajectory) rules out moving-target hypothesis.]

\subsection{Component Ablations}
\label{sec:results:component_ablations}

% Source: Awaiting ablation results (Day 14-16).

[To be added Day~14--16 once ablation results are available, via
\texttt{scripts/diagnostic/25\_v\_d\_ablation\_aggregate.py}: 21-cell ablation design
(Path~B base, Path~B + Cocos, variant~(c), V3 baseline, drop-mesh, drop-grasp,
drop-proprio; each $\times$ 3 seeds). Statistical protocol per preregistration: pairwise
McNemar + BH-FDR adjustment, Wilson 95\% CI per cell--task, Welch's $t$-test on per-seed
SR.]
